\section{Application}

The scheme was evaluated only on simulated data; no industrial or benchmark process, such as the
Tennessee Eastman Process, was monitored in this study. In practice it would be deployed in the same
two phases. In Phase~I, a stable reference period of normal operation is divided chronologically into
three parts that mirror D1--D3: standardisation and CAE training, OC-SVM fitting on the frozen
encoder's features, and estimation of the moments of $y_t$ with empirical calibration of $h^*$ to the
in-control ARL the plant can tolerate. In Phase~II, each new observation completes a window of the 32
most recent standardised observations; one CAE pass and one OC-SVM evaluation give $s_t$ and
$\ell_t$, which update the MEWMA, and a signal prompts investigation rather than automatic correction.
The computation is light: the CAE has 3180 parameters and the OC-SVM 78 support vectors, and the
complete production run of this study (data generation, training, calibration and Phase~II
evaluation of all six charts) took 207 seconds on a single laptop processor thread.

The practical demands are nonetheless real. The reference period must be genuinely in control,
because faults present in it would be learned as normal. The persistent window features make the
first reports after a restart conservative: no in-control D4 trajectory of $H(0.05)$ signalled before
its eleventh report, whereas $M(0.05)$ could signal at its first. Legitimate process changes require
refitting and recalibration on fresh in-control data, and a signal is harder to trace to individual
variables than with a raw-data MEWMA, whose state vector indicates the direction of a shift
\citep{lowry_multivariate_1992}. Above all, the evidence concerns one simulated process and one fault
family, for which the raw-data benchmark was faster; external validation on the Tennessee Eastman
Process or real plant data is needed before industrial use can be recommended.
